\documentclass[12pt,a4paper]{article}

\usepackage[utf8]{inputenc}
\usepackage[T1]{fontenc}
\usepackage{amsmath,amssymb}
\usepackage{geometry}
\usepackage{booktabs}
\usepackage{array}
\usepackage{xcolor}
\usepackage{enumerate}

\geometry{margin=2.3cm}
\setlength{\parskip}{0.5em}
\setlength{\parindent}{0pt}
\setlength{\jot}{8pt}
\pagestyle{plain}

\begin{document}

% ===================================================================
% ENCABEZADO
% ===================================================================
\begin{center}
    {\LARGE \textbf{PRÁCTICA POR EQUIPO}}\\[0.3cm]
    {\Large \textbf{Procesos Markovianos de Decisión}}\\[0.2cm]
    {\large \textbf{Optimización de Estrategia Publicitaria}}\\[0.5cm]
    \rule{0.8\textwidth}{0.4pt}\\[0.3cm]
    \textbf{Integrantes:} \_\_\_\_\_\_\_\_\_\_\_\_\_\_\_\_\_\_\_\_\_\_\_\_\_\_\_\_\_\_\_\_\_\_\_\_\_\_\_\_\_\_\_\_\_\_\\[0.2cm]
    \textbf{Fecha:} \_\_\_\_\_\_\_\_\_\_\_\_\_\_\_\_\_\_\_\_\_\_\_\_\_\_\_\_\_\_\\[0.3cm]
    \rule{0.8\textwidth}{0.4pt}
\end{center}

\vspace{0.3cm}

\section*{Instrucciones}

\begin{itemize}
    \item Resolver en hojas limpias, mostrando \textbf{todo} el desarrollo.
    \item Si se necesita factor de descuento, usar $\alpha = 0.9$.
    \item Si se necesita política de arranque, usar $R_0 = [1,\; 1,\; 2]$: \textbf{Radio} en Regular, \textbf{Radio} en Bueno, \textbf{TV} en Excelente.
    \item En cualquier método iterativo, elaborar un \textbf{máximo de 2 iteraciones}.
\end{itemize}

\newpage

% ===================================================================
% 1. ENUNCIADO
% ===================================================================
\section{Enunciado}

Una empresa elige semanalmente entre tres medios publicitarios para \textbf{maximizar} utilidades. Las ventas se clasifican en: \textbf{Regular} (0), \textbf{Bueno} (1) y \textbf{Excelente} (2).

\begin{center}
\begin{tabular}{c c c}
\toprule
\textbf{Decisión} & \textbf{Medio} & \textbf{Costo fijo} \\
\midrule
$k = 1$ & Radio      & \$200 \\
$k = 2$ & TV         & \$900 \\
$k = 3$ & Periódico  & \$300 \\
\bottomrule
\end{tabular}
\end{center}

Distribución inicial: $a = [0.4,\; 0.3,\; 0.3]$.

\vspace{0.3cm}
\textbf{Se pide:}

\begin{enumerate}[\textbf{A)}]
    \item Definir el espacio de estados $E$, decisiones $K$, y costos fijos.
    \item Escribir las matrices de transición $P^{(k)}$ y de ingreso $Ingreso^{(k)}$.
    \item Calcular \textit{a mano} la matriz de costos $C_{ik}$ con:
    \[ C_{ik}^{\text{int}} = -\Bigl( \sum_{j} P_{ij}(k) \cdot Ingreso_{ij}(k) \;-\; CostoFijo(k) \Bigr) \]
    (El signo negativo convierte maximización en minimización.)
    \item ¿Cuántas políticas determinísticas hay? Enumerarlas.
    \item \textbf{Mejoramiento sin descuento:} 2 iteraciones desde $R_0 = [1,1,2]$.
    \item \textbf{Mejoramiento con descuento} ($\alpha = 0.9$): 2 iteraciones desde $R_0 = [1,1,2]$.
    \item Plantear el Problema de Programación Lineal asociado.
\end{enumerate}

\newpage

% ===================================================================
% A) ESTADOS, DECISIONES, COSTOS
% ===================================================================
\section{Inciso A --- Estados, decisiones y costos}

\subsection*{Espacio de estados}

\[ E = \{0,\; 1,\; 2\}, \qquad m = 2 \]

\begin{center}
\begin{tabular}{c c l}
\toprule
\textbf{Índice} & \textbf{Estado} & \textbf{Significado} \\
\midrule
0 & Regular   & Ventas semanales bajas \\
1 & Bueno     & Ventas semanales medias \\
2 & Excelente & Ventas semanales altas \\
\bottomrule
\end{tabular}
\end{center}

\subsection*{Conjunto de decisiones}

$K = 3$ medios disponibles:

\begin{center}
\begin{tabular}{c c c c}
\toprule
\textbf{Notación} & \textbf{Decisión} & \textbf{Medio} & \textbf{Costo fijo}\\
\midrule
$k = 1$ & 1 & Radio     & \$200 \\
$k = 2$ & 2 & TV        & \$900 \\
$k = 3$ & 3 & Periódico & \$300 \\
\bottomrule
\end{tabular}
\end{center}

\subsection*{Datos adicionales}

\begin{itemize}
    \item \textbf{Vector inicial:} $a = [0.4,\; 0.3,\; 0.3]$.
    \item \textbf{Objetivo:} maximizar $\Rightarrow$ los costos se niegan: $\min(-\text{utilidad}) \equiv \max(\text{utilidad})$.
    \item \textbf{Factor de descuento} (cuando aplique): $\alpha = 0.9$.
\end{itemize}

\newpage

% ===================================================================
% B) MATRICES DE TRANSICIÓN E INGRESO
% ===================================================================
\section{Inciso B --- Matrices de transición y de ingreso}

Cada decisión $k$ define una matriz $P^{(k)}$ (probabilidades de transición) y una matriz $Ingreso^{(k)}$ (ganancia bruta por transición, antes de restar el costo fijo).

\subsection*{Radio ($k = 1$) --- Costo fijo: \$200}

\[
P^{(1)} = \begin{bmatrix}
0.4 & 0.5 & 0.1 \\
0.1 & 0.7 & 0.2 \\
0.1 & 0.2 & 0.7
\end{bmatrix}
\qquad
Ingreso^{(1)} = \begin{bmatrix}
400 & 520 & 600 \\
300 & 400 & 700 \\
200 & 250 & 500
\end{bmatrix}
\]

\subsection*{TV ($k = 2$) --- Costo fijo: \$900}

\[
P^{(2)} = \begin{bmatrix}
0.7 & 0.2 & 0.1 \\
0.3 & 0.6 & 0.1 \\
0.1 & 0.7 & 0.2
\end{bmatrix}
\qquad
Ingreso^{(2)} = \begin{bmatrix}
1000 & 1300 & 1600 \\
800  & 1000 & 1700 \\
600  & 700  & 1100
\end{bmatrix}
\]

\subsection*{Periódico ($k = 3$) --- Costo fijo: \$300}

\[
P^{(3)} = \begin{bmatrix}
0.2 & 0.5 & 0.3 \\
0.0 & 0.7 & 0.3 \\
0.0 & 0.2 & 0.8
\end{bmatrix}
\qquad
Ingreso^{(3)} = \begin{bmatrix}
400 & 530 & 710 \\
350 & 450 & 800 \\
250 & 400 & 650
\end{bmatrix}
\]

\vspace{0.3cm}
\begin{quote}
\textbf{Observación importante:} En $P^{(3)}$ (Periódico), $P_{10}(3)=0$ y $P_{20}(3)=0$. Desde Bueno y Excelente \textbf{nunca} se regresa a Regular. Esto hará transitorio al estado 0 bajo la política óptima.
\end{quote}

\newpage

% ===================================================================
% C) CÁLCULO DE C_ik
% ===================================================================
\section{Inciso C --- Cálculo manual de la matriz de costos $C_{ik}$}

\subsection*{Fórmula general}

\[
C_{ik}^{\,\text{int}} = -\Bigl(\; \underbrace{\sum_{j=0}^{2} P_{ij}(k) \cdot Ingreso_{ij}(k)}_{\text{Ingreso esperado}} \;-\; CostoFijo(k) \Bigr)
\]

La \textbf{utilidad neta real} es el valor sin negar. El signo negativo convierte maximización en minimización.

\subsection*{Radio ($k = 1$), costo fijo = \$200}

\vspace{0.2cm}
\textbf{Estado 0 (Regular) ---} Fila 0 de $P^{(1)}$ y de $Ingreso^{(1)}$:

\[
\begin{aligned}
\mathbb{E}[\text{Ingreso}] &= 0.4(400) + 0.5(520) + 0.1(600) \\
&= 160 + 260 + 60 = 480
\end{aligned}
\]

\[
\text{Utilidad neta} = 480 - 200 = 280
\qquad\Longrightarrow\qquad
\boxed{C_{0,1} = -280}
\]

\vspace{0.3cm}
\textbf{Estado 1 (Bueno) ---} Fila 1:

\[
\begin{aligned}
\mathbb{E}[\text{Ingreso}] &= 0.1(300) + 0.7(400) + 0.2(700) \\
&= 30 + 280 + 140 = 450
\end{aligned}
\qquad\Longrightarrow\qquad
\boxed{C_{1,1} = -250}
\]

\vspace{0.3cm}
\textbf{Estado 2 (Excelente) ---} Fila 2:

\[
\begin{aligned}
\mathbb{E}[\text{Ingreso}] &= 0.1(200) + 0.2(250) + 0.7(500) \\
&= 20 + 50 + 350 = 420
\end{aligned}
\qquad\Longrightarrow\qquad
\boxed{C_{2,1} = -220}
\]

\subsection*{TV ($k = 2$), costo fijo = \$900}

\textbf{Estado 0:}
\[
\mathbb{E}[\text{Ingreso}] = 0.7(1000) + 0.2(1300) + 0.1(1600) = 700+260+160 = 1120
\quad\Longrightarrow\quad
\boxed{C_{0,2} = -220}
\]

\textbf{Estado 1:}
\[
\mathbb{E}[\text{Ingreso}] = 0.3(800) + 0.6(1000) + 0.1(1700) = 240+600+170 = 1010
\quad\Longrightarrow\quad
\boxed{C_{1,2} = -110}
\]

\textbf{Estado 2:}
\[
\mathbb{E}[\text{Ingreso}] = 0.1(600) + 0.7(700) + 0.2(1100) = 60+490+220 = 770
\]

\[
\text{Utilidad neta} = 770 - 900 = -130 \;\; \text{\color{red}(\textbf{¡pérdida!})}
\quad\Longrightarrow\quad
\boxed{C_{2,2} = -(-130) = \mathbf{+130}}
\]

\textbf{Interpretación:} TV en Excelente es la única combinación con pérdida neta (-\$130). Su alto costo fijo (\$900) supera el ingreso esperado (\$770). En la matriz de costos aparece como \textbf{+130} --- el único valor positivo ---, que los algoritmos de minimización evitarán.

\subsection*{Periódico ($k = 3$), costo fijo = \$300}

\textbf{Estado 0:}
\[
\mathbb{E}[\text{Ingreso}] = 0.2(400) + 0.5(530) + 0.3(710) = 80+265+213 = 558
\quad\Longrightarrow\quad
\boxed{C_{0,3} = -258}
\]

\textbf{Estado 1:}
\[
\mathbb{E}[\text{Ingreso}] = 0.0(350) + 0.7(450) + 0.3(800) = 0+315+240 = 555
\quad\Longrightarrow\quad
\boxed{C_{1,3} = -255}
\]

\textbf{Estado 2:}
\[
\mathbb{E}[\text{Ingreso}] = 0.0(250) + 0.2(400) + 0.8(650) = 0+80+520 = 600
\quad\Longrightarrow\quad
\boxed{C_{2,3} = -300}
\]

\subsection*{Matriz de costos resultante}

\[
\boxed{ C = \begin{bmatrix}
-280 & -220 & -258 \\
-250 & -110 & -255 \\
-220 & \mathbf{+130} & -300
\end{bmatrix} }
\qquad
\text{Utilidad neta real } U = -C = \begin{bmatrix}
280 & 220 & 258 \\
250 & 110 & 255 \\
220 & \mathbf{-130} & 300
\end{bmatrix}
\]

\newpage

% ===================================================================
% D) POLÍTICAS
% ===================================================================
\section{Inciso D --- Políticas determinísticas}

Una \textbf{política} $R = [d(0),\, d(1),\, d(2)]$ asigna una decisión a cada estado: $d(i) \in \{1,2,3\}$.

\[
\text{Total: } K^{\,m+1} = 3^{3} = \mathbf{27 \text{ políticas}}
\]

\begin{center}
\footnotesize
\begin{tabular}{c c c | c c c | c c c}
\toprule
$[1,1,1]$ & $[1,1,2]$ & $[1,1,3]$ & $[1,2,1]$ & $[1,2,2]$ & $[1,2,3]$ & $[1,3,1]$ & $[1,3,2]$ & $[1,3,3]$ \\
$[2,1,1]$ & $[2,1,2]$ & $[2,1,3]$ & $[2,2,1]$ & $[2,2,2]$ & $[2,2,3]$ & $[2,3,1]$ & $[2,3,2]$ & $[2,3,3]$ \\
$[3,1,1]$ & $[3,1,2]$ & $[3,1,3]$ & $[3,2,1]$ & $[3,2,2]$ & $[3,2,3]$ & $[3,3,1]$ & $[3,3,2]$ & $[3,3,3]$ \\
\bottomrule
\end{tabular}
\end{center}

\vspace{0.3cm}
Política de arranque: $R_0 = [1, 1, 2]$, es decir $(\text{Radio},\, \text{Radio},\, \text{TV})$.

\newpage

% ===================================================================
% E) MEJORAMIENTO SIN DESCUENTO
% ===================================================================
\section{Inciso E --- Mejoramiento de políticas sin descuento}

\textbf{Algoritmo de Howard para costo promedio.} Consta de dos pasos:
\begin{enumerate}
    \item[\textbf{Paso 1.}] \textbf{Determinación del valor:} se calcula el costo promedio $g$ y los valores relativos $V_i$ de cada estado bajo la política actual.
    \item[\textbf{Paso 2.}] \textbf{Mejoramiento:} para cada estado se busca si otra decisión da un menor costo ajustado. Si hay cambios, se repite desde el Paso 1.
\end{enumerate}

\vspace{0.2cm}
\textbf{Política de arranque:} $R_0 = [1, 1, 2] = (\text{Radio}, \text{Radio}, \text{TV})$.

% -------------------------------------------------------------------
\subsection*{Iteración 1}

\subsubsection*{Paso 1 --- Determinación del valor}

Construimos la matriz de transición combinada $P_{R_0}$: para cada estado $i$, se toma la fila $i$ de la matriz $P^{(d(i))}$ correspondiente a la decisión indicada por $R_0$:
\begin{itemize}
    \item Estado 0 $\to$ decisión 1 (Radio) $\to$ fila 0 de $P^{(1)}$: $[0.4,\; 0.5,\; 0.1]$
    \item Estado 1 $\to$ decisión 1 (Radio) $\to$ fila 1 de $P^{(1)}$: $[0.1,\; 0.7,\; 0.2]$
    \item Estado 2 $\to$ decisión 2 (TV) $\to$ fila 2 de $P^{(2)}$: $[0.1,\; 0.7,\; 0.2]$
\end{itemize}

\[
P_{R_0} = \begin{bmatrix}
0.4 & 0.5 & 0.1 \\
0.1 & 0.7 & 0.2 \\
0.1 & 0.7 & 0.2
\end{bmatrix}
\qquad
C_{R_0} = \begin{bmatrix} C_{0,1} \\ C_{1,1} \\ C_{2,2} \end{bmatrix}
= \begin{bmatrix} -280 \\ -250 \\ +130 \end{bmatrix}
\]

Planteamos el sistema $g + V_i - \sum_{j} P_{ij}^{R_0}\, V_j = C_{i,R_0(i)}$, con $V_2 = 0$:

\[
\begin{cases}
g + V_0 - (0.4V_0 + 0.5V_1 + 0.1\cdot 0) = -280 \\
g + V_1 - (0.1V_0 + 0.7V_1 + 0.2\cdot 0) = -250 \\
g + 0 \; - (0.1V_0 + 0.7V_1 + 0.2\cdot 0) = 130
\end{cases}
\]

Simplificando:

\[
\boxed{\begin{cases}
g + 0.6V_0 - 0.5V_1 = -280 & \text{(E1)}\\[4pt]
g - 0.1V_0 + 0.3V_1 = -250 & \text{(E2)}\\[4pt]
g - 0.1V_0 - 0.7V_1 = 130  & \text{(E3)}
\end{cases}}
\]

\textbf{Resolución:} De (E3): $g = 130 + 0.1V_0 + 0.7V_1$. Sustituyendo en (E2):

\[
(130 + 0.1V_0 + 0.7V_1) - 0.1V_0 + 0.3V_1 = -250
\;\Longrightarrow\; 130 + V_1 = -250
\;\Longrightarrow\; \boxed{V_1 = -380}
\]

Ahora $g = 130 + 0.1V_0 + 0.7(-380) = -136 + 0.1V_0$. Sustituyendo en (E1):

\[
(-136 + 0.1V_0) + 0.6V_0 - 0.5(-380) = -280
\;\Longrightarrow\; 0.7V_0 + 54 = -280
\;\Longrightarrow\; \boxed{V_0 = -\tfrac{334}{0.7} \approx -477.14}
\]

\[
g = -136 + 0.1(-477.14) \approx \boxed{-183.71}
\]

\[
\boxed{ \boxed{ g = -183.71,\quad V_0 = -477.14,\quad V_1 = -380,\quad V_2 = 0 } }
\]

\textbf{Interpretación:} La utilidad promedio actual es \$183.71/semana. Los valores $V_i$ indican qué tan ``costoso'' es cada estado con respecto al estado 2 (Excelente).

\subsubsection*{Paso 2 --- Mejoramiento}

Para cada estado $i$ y decisión $k$, calculamos:
\[ \text{test}(i,k) = C_{ik} + \sum_{j=0}^{2} P_{ij}(k)\, V_j \;-\; V_i \]
La decisión con el \textbf{menor} test reemplaza a la actual.

\vspace{0.3cm}
\textbf{Estado 0 (Regular), $V_0 = -477.14$:}

\begin{quote}\footnotesize
$k=1$ (Radio):\;
$\sum P_{0j}V_j = 0.4(-477.14) + 0.5(-380) + 0.1(0) = -190.86 - 190.00 = -380.86$ \\
\hspace*{1.8cm} $\text{test} = C_{0,1} + \sum P_{0j}V_j - V_0 = -280 + (-380.86) - (-477.14) = \mathbf{-183.71}$\quad $\leftarrow$ \textbf{MEJOR}

\vspace{4pt}
$k=2$ (TV):\;
$\sum P_{0j}V_j = 0.7(-477.14) + 0.2(-380) + 0.1(0) = -334.00 - 76.00 = -410.00$ \\
\hspace*{1.8cm} $\text{test} = -220 + (-410.00) - (-477.14) = -152.86$

\vspace{4pt}
$k=3$ (Periódico):\;
$\sum P_{0j}V_j = 0.2(-477.14) + 0.5(-380) + 0.3(0) = -95.43 - 190.00 = -285.43$ \\
\hspace*{1.8cm} $\text{test} = -258 + (-285.43) - (-477.14) = -66.29$
\end{quote}

\vspace{0.3cm}
\textbf{Estado 1 (Bueno), $V_1 = -380$:}

\begin{quote}\footnotesize
$k=1$ (Radio):\;
$\sum P_{1j}V_j = 0.1(-477.14) + 0.7(-380) + 0.2(0) = -47.71 - 266.00 = -313.71$ \\
\hspace*{1.8cm} $\text{test} = -250 + (-313.71) - (-380) = \mathbf{-183.71}$\quad $\leftarrow$ \textbf{MEJOR}

\vspace{4pt}
$k=2$ (TV):\;
$\sum P_{1j}V_j = 0.3(-477.14) + 0.6(-380) + 0.1(0) = -143.14 - 228.00 = -371.14$ \\
\hspace*{1.8cm} $\text{test} = -110 + (-371.14) - (-380) = -101.14$

\vspace{4pt}
$k=3$ (Periódico):\;
$\sum P_{1j}V_j = 0.0(-477.14) + 0.7(-380) + 0.3(0) = -266.00$ \\
\hspace*{1.8cm} $\text{test} = -255 + (-266.00) - (-380) = -141.00$
\end{quote}

\vspace{0.3cm}
\textbf{Estado 2 (Excelente), $V_2 = 0$ (no se resta $V_i$):}

\begin{quote}\footnotesize
$k=1$ (Radio):\;
$\sum P_{2j}V_j = 0.1(-477.14) + 0.2(-380) + 0.7(0) = -47.71 - 76.00 = -123.71$ \\
\hspace*{1.8cm} $\text{test} = -220 + (-123.71) = -343.71$

\vspace{4pt}
$k=2$ (TV):\;
$\sum P_{2j}V_j = 0.1(-477.14) + 0.7(-380) + 0.2(0) = -47.71 - 266.00 = -313.71$ \\
\hspace*{1.8cm} $\text{test} = +130 + (-313.71) = -183.71$

\vspace{4pt}
$k=3$ (Periódico):\;
$\sum P_{2j}V_j = 0.0(-477.14) + 0.2(-380) + 0.8(0) = -76.00$ \\
\hspace*{1.8cm} $\text{test} = -300 + (-76.00) = \mathbf{-376.00}$\quad $\leftarrow$ \textbf{MEJOR}
\end{quote}

\vspace{0.3cm}
\textbf{Resultado:} Solo cambia el estado 2 (TV $\to$ Periódico). TV tenía $C_{2,2}=+130$ (pérdida), lo que la hace muy desfavorable frente a Periódico ($C_{2,3}=-300$).

\[
\boxed{ R_1 = [1,\; 1,\; 3] = (\text{Radio},\; \text{Radio},\; \text{Periódico}) }
\]

\newpage

% -------------------------------------------------------------------
\subsection*{Iteración 2}

Ahora $R_1 = [1, 1, 3]$.

\subsubsection*{Paso 1 --- Determinación del valor}

Matriz combinada (Radio en 0,1; Periódico en 2):

\[
P_{R_1} = \begin{bmatrix}
0.4 & 0.5 & 0.1 \\
0.1 & 0.7 & 0.2 \\
0.0 & 0.2 & 0.8
\end{bmatrix}
\qquad
C_{R_1} = \begin{bmatrix} -280 \\ -250 \\ -300 \end{bmatrix}
\]

Nótese el cambio en la tercera fila: $P_{20}=0$ (con Periódico no se cae a Regular desde Excelente).

Sistema con $V_2 = 0$:

\[
\boxed{\begin{cases}
g + 0.6V_0 - 0.5V_1 = -280 & \text{(E1)}\\[4pt]
g - 0.1V_0 + 0.3V_1 = -250 & \text{(E2)}\\[4pt]
g - 0.2V_1 = -300 & \text{(E3)}
\end{cases}}
\]

De (E3): $g = -300 + 0.2V_1$. En (E2):
\[
(-300 + 0.2V_1) - 0.1V_0 + 0.3V_1 = -250
\;\Longrightarrow\; -0.1V_0 + 0.5V_1 = 50
\;\Longrightarrow\; \boxed{V_1 = 100 + 0.2V_0}
\]

Entonces $g = -300 + 0.2(100 + 0.2V_0) = -280 + 0.04V_0$. En (E1):
\[
(-280 + 0.04V_0) + 0.6V_0 - 0.5(100 + 0.2V_0) = -280
\;\Longrightarrow\; 0.54V_0 - 330 = -280
\;\Longrightarrow\; \boxed{V_0 = \tfrac{50}{0.54} \approx 92.59}
\]

\[
V_1 = 100 + 0.2(92.59) \approx 118.52,\qquad
g = -280 + 0.04(92.59) \approx -276.30
\]

\[
\boxed{ \boxed{ g = -276.30,\quad V_0 = 92.59,\quad V_1 = 118.52,\quad V_2 = 0 } }
\]

\textbf{Cambio notable:} La utilidad subió de \$183.71 a \$276.30/semana. Los $V_i$ ahora son \textbf{positivos} (salvo $V_2=0$): estar en Regular o Bueno es ``mejor'' que estar en Excelente bajo esta política.

\subsubsection*{Paso 2 --- Mejoramiento}

\textbf{Estado 0 (Regular), $V_0 = 92.59$:}

\begin{quote}\footnotesize
$k=1$ (Radio):\;
$\sum P_{0j}V_j = 0.4(92.59) + 0.5(118.52) + 0.1(0) = 37.04 + 59.26 = 96.30$ \\
\hspace*{1.8cm} $\text{test} = -280 + 96.30 - 92.59 = \mathbf{-276.29}$\quad $\leftarrow$ \textbf{MEJOR}

\vspace{4pt}
$k=2$ (TV):\;
$\sum P_{0j}V_j = 0.7(92.59) + 0.2(118.52) + 0.1(0) = 64.81 + 23.70 = 88.52$ \\
\hspace*{1.8cm} $\text{test} = -220 + 88.52 - 92.59 = -224.07$

\vspace{4pt}
$k=3$ (Periódico):\;
$\sum P_{0j}V_j = 0.2(92.59) + 0.5(118.52) + 0.3(0) = 18.52 + 59.26 = 77.78$ \\
\hspace*{1.8cm} $\text{test} = -258 + 77.78 - 92.59 = -272.81$
\end{quote}

\vspace{0.3cm}
\textbf{Estado 1 (Bueno), $V_1 = 118.52$:}

\begin{quote}\footnotesize
$k=1$ (Radio):\;
$\sum P_{1j}V_j = 0.1(92.59) + 0.7(118.52) + 0.2(0) = 9.26 + 82.96 = 92.22$ \\
\hspace*{1.8cm} $\text{test} = -250 + 92.22 - 118.52 = -276.30$

\vspace{4pt}
$k=2$ (TV):\;
$\sum P_{1j}V_j = 0.3(92.59) + 0.6(118.52) + 0.1(0) = 27.78 + 71.11 = 98.89$ \\
\hspace*{1.8cm} $\text{test} = -110 + 98.89 - 118.52 = -129.63$

\vspace{4pt}
$k=3$ (Periódico):\;
$\sum P_{1j}V_j = 0.0(92.59) + 0.7(118.52) + 0.3(0) = 82.96$ \\
\hspace*{1.8cm} $\text{test} = -255 + 82.96 - 118.52 = \mathbf{-290.56}$\quad $\leftarrow$ \textbf{MEJOR}
\end{quote}

\textbf{Explicación:} Aunque Radio es más barato (\$200 vs \$300), Periódico gana en el estado Bueno porque su test ($-290.56$) es menor (más negativo) que el de Radio ($-276.30$). La razón: Periódico tiene $P_{10}(3)=0$, nunca degrada a Regular, y su ingreso esperado es mayor.

\vspace{0.3cm}
\textbf{Estado 2 (Excelente), $V_2 = 0$:}

\begin{quote}\footnotesize
$k=1$ (Radio):\;
$\sum P_{2j}V_j = 0.1(92.59) + 0.2(118.52) + 0.7(0) = 9.26 + 23.70 = 32.96$ \\
\hspace*{1.8cm} $\text{test} = -220 + 32.96 = -187.04$

\vspace{4pt}
$k=2$ (TV):\;
$\sum P_{2j}V_j = 0.1(92.59) + 0.7(118.52) + 0.2(0) = 9.26 + 82.96 = 92.22$ \\
\hspace*{1.8cm} $\text{test} = +130 + 92.22 = \mathbf{+222.22}$\quad (¡único test positivo!)

\vspace{4pt}
$k=3$ (Periódico):\;
$\sum P_{2j}V_j = 0.0(92.59) + 0.2(118.52) + 0.8(0) = 23.70$ \\
\hspace*{1.8cm} $\text{test} = -300 + 23.70 = \mathbf{-276.30}$\quad $\leftarrow$ \textbf{MEJOR}
\end{quote}

\textbf{Observación sobre TV:} En Excelente, TV da test $=+222.22$ --- el único positivo de toda la tabla. Esto confirma que TV es pésima elección allí.

\vspace{0.3cm}
\textbf{Resultado iteración 2:} Cambia estado 1 de Radio a Periódico.

\[
\boxed{ R_2 = [1,\; 3,\; 3] = (\text{Radio},\; \text{Periódico},\; \text{Periódico}) }
\]

En solo 2 iteraciones se alcanzó $[1,3,3]$, la \textbf{política óptima}.

\newpage

% ===================================================================
% F) MEJORAMIENTO CON DESCUENTO
% ===================================================================
\section{Inciso F --- Mejoramiento con descuento ($\alpha = 0.9$)}

Con descuento, los valores $V_i$ representan el \textbf{costo total descontado} a largo plazo empezando en cada estado. La ecuación de evaluación es:
\[ V_i = C_{i,R(i)} + \alpha \sum_{j} P_{ij}(R(i))\, V_j \]
Y el criterio de mejoramiento es $\text{valor}(i,k) = C_{ik} + \alpha\sum_{j}P_{ij}(k)V_j$ (sin restar $V_i$).

\vspace{0.2cm}
\textbf{Política de arranque:} $R_0 = [1, 1, 2] = (\text{Radio}, \text{Radio}, \text{TV})$.

% -------------------------------------------------------------------
\subsection*{Iteración 1}

\subsubsection*{Paso 1 --- Determinación del valor con descuento}

Usando $P_{R_0}$ y $\alpha = 0.9$:

\[
\begin{cases}
V_0 = -280 + 0.9(0.4V_0 + 0.5V_1 + 0.1V_2) \\
V_1 = -250 + 0.9(0.1V_0 + 0.7V_1 + 0.2V_2) \\
V_2 = 130 \;+\; 0.9(0.1V_0 + 0.7V_1 + 0.2V_2)
\end{cases}
\]

Reordenando:

\[
\boxed{\begin{cases}
0.64V_0 - 0.45V_1 - 0.09V_2 = -280 & \text{(D1)}\\[4pt]
-0.09V_0 + 0.37V_1 - 0.18V_2 = -250 & \text{(D2)}\\[4pt]
-0.09V_0 - 0.63V_1 + 0.82V_2 = 130 & \text{(D3)}
\end{cases}}
\]

\textbf{Resolución.} De (D3):
\[
V_2 = \frac{130}{0.82} + \frac{0.09}{0.82}V_0 + \frac{0.63}{0.82}V_1
     = 158.54 + 0.1098\,V_0 + 0.7683\,V_1
\]

Sustituyendo $V_2$ en (D2):
\[
\begin{aligned}
-0.09V_0 + 0.37V_1 - 0.18(158.54 + 0.1098V_0 + 0.7683V_1) &= -250 \\
-0.1098V_0 + 0.2317V_1 - 28.54 &= -250 \\
0.2317V_1 &= 0.1098V_0 - 221.46
\end{aligned}
\]
\[
\boxed{V_1 = 0.4739\,V_0 - 955.81}
\]

Ahora actualizamos $V_2$ con esta expresión de $V_1$:
\[
\begin{aligned}
V_2 &= 158.54 + 0.1098V_0 + 0.7683(0.4739V_0 - 955.81) \\
    &= 158.54 + 0.1098V_0 + 0.3640V_0 - 734.47 \\
    &= -575.93 + 0.4738\,V_0
\end{aligned}
\]

Sustituyendo $V_1$ y $V_2$ en (D1):
\[
\begin{aligned}
0.64V_0 - 0.45(0.4739V_0 - 955.81) - 0.09(-575.93 + 0.4738V_0) &= -280 \\
0.64V_0 - 0.2133V_0 + 430.11 + 51.83 - 0.0426V_0 &= -280 \\
(0.64 - 0.2133 - 0.0426)V_0 + 481.94 &= -280 \\
0.3841V_0 &= -761.94
\end{aligned}
\]

\[
\boxed{V_0 \approx -1983.77}
\]

\[
V_1 = 0.4739(-1983.77) - 955.81 \approx -940.12 - 955.81 \approx \boxed{-1895.93}
\]

\[
V_2 = -575.93 + 0.4738(-1983.77) \approx -575.93 - 939.96 \approx \boxed{-1515.89}
\]

\[
\boxed{ \boxed{ V_0 = -1983.77,\quad V_1 = -1895.93,\quad V_2 = -1515.89 } }
\]

\textbf{Interpretación:} Los $V_i$ son la utilidad total descontada esperada a largo plazo partiendo de cada estado. Los tres valores son negativos (utilidad). El estado Excelente tiene el ``menos negativo'' ($-1516$) porque bajo $R_0$ se usa TV allí, que genera pérdida ($C_{2,2}=+130$). El estado Regular ($-1984$) acumula más utilidad porque Radio es más rentable.

\subsubsection*{Paso 2 --- Mejoramiento con descuento}

El criterio es $\text{valor}(i,k) = C_{ik} + \alpha\sum_{j} P_{ij}(k)V_j$. Se elige el \textbf{menor} valor. Nótese que \textbf{no se resta $V_i$} a diferencia del caso sin descuento.

\vspace{0.3cm}
\textbf{Estado 0:}

\begin{quote}\footnotesize
$k=1$ (Radio):\;
$\alpha\sum P_{0j}V_j = 0.9[0.4(-1983.77) + 0.5(-1895.93) + 0.1(-1515.89)]$ \\
\hspace*{1.8cm} $= 0.9[-793.5 - 948.0 - 151.6] = -1703.8$ \\
\hspace*{1.8cm} $\text{valor} = -280 + (-1703.8) = \mathbf{-1983.8}$\quad $\leftarrow$ \textbf{MEJOR}

\vspace{3pt}
$k=2$ (TV):\;
$\alpha\sum P_{0j}V_j = 0.9[0.7(-1983.77) + 0.2(-1895.93) + 0.1(-1515.89)]$ \\
\hspace*{1.8cm} $= 0.9[-1388.6 - 379.2 - 151.6] = -1727.5$ \\
\hspace*{1.8cm} $\text{valor} = -220 + (-1727.5) = -1947.5$

\vspace{3pt}
$k=3$ (Periódico):\;
$\alpha\sum P_{0j}V_j = 0.9[0.2(-1983.77) + 0.5(-1895.93) + 0.3(-1515.89)]$ \\
\hspace*{1.8cm} $= 0.9[-396.8 - 948.0 - 454.8] = -1619.6$ \\
\hspace*{1.8cm} $\text{valor} = -258 + (-1619.6) = -1877.6$
\end{quote}

\vspace{0.3cm}
\textbf{Estado 1:}

\begin{quote}\footnotesize
$k=1$ (Radio):\;
$\alpha\sum P_{1j}V_j = 0.9[0.1(-1983.77) + 0.7(-1895.93) + 0.2(-1515.89)]$ \\
\hspace*{1.8cm} $= 0.9[-198.4 - 1327.2 - 303.2] = -1645.9$ \\
\hspace*{1.8cm} $\text{valor} = -250 + (-1645.9) = \mathbf{-1895.9}$\quad $\leftarrow$ \textbf{MEJOR}

\vspace{3pt}
$k=2$ (TV):\;
$\alpha\sum P_{1j}V_j = 0.9[0.3(-1983.77) + 0.6(-1895.93) + 0.1(-1515.89)]$ \\
\hspace*{1.8cm} $= 0.9[-595.1 - 1137.6 - 303.2] = -1832.3$ \\
\hspace*{1.8cm} $\text{valor} = -110 + (-1832.3) = -1942.3$

\vspace{3pt}
$k=3$ (Periódico):\;
$\alpha\sum P_{1j}V_j = 0.9[0.0(-1983.77) + 0.7(-1895.93) + 0.3(-1515.89)]$ \\
\hspace*{1.8cm} $= 0.9[0 - 1327.2 - 454.8] = -1603.7$ \\
\hspace*{1.8cm} $\text{valor} = -255 + (-1603.7) = -1858.7$
\end{quote}

\vspace{0.3cm}
\textbf{Estado 2:}

\begin{quote}\footnotesize
$k=1$ (Radio):\;
$\alpha\sum P_{2j}V_j = 0.9[0.1(-1983.77) + 0.2(-1895.93) + 0.7(-1515.89)]$ \\
\hspace*{1.8cm} $= 0.9[-198.4 - 379.2 - 1061.1] = -1474.8$ \\
\hspace*{1.8cm} $\text{valor} = -220 + (-1474.8) = -1694.8$

\vspace{3pt}
$k=2$ (TV):\;
$\alpha\sum P_{2j}V_j = 0.9[0.1(-1983.77) + 0.7(-1895.93) + 0.2(-1515.89)]$ \\
\hspace*{1.8cm} $= 0.9[-198.4 - 1327.2 - 303.2] = -1645.9$ \\
\hspace*{1.8cm} $\text{valor} = +130 + (-1645.9) = -1515.9$

\vspace{3pt}
$k=3$ (Periódico):\;
$\alpha\sum P_{2j}V_j = 0.9[0.0(-1983.77) + 0.2(-1895.93) + 0.8(-1515.89)]$ \\
\hspace*{1.8cm} $= 0.9[0 - 379.2 - 1212.7] = -1432.7$ \\
\hspace*{1.8cm} $\text{valor} = -300 + (-1432.7) = \mathbf{-1732.7}$\quad $\leftarrow$ \textbf{MEJOR}
\end{quote}

\[
\boxed{ R_1 = [1,\; 1,\; 3] = (\text{Radio},\; \text{Radio},\; \text{Periódico}) }
\]

Mismo cambio que en el caso sin descuento: estado 2 cambia de TV a Periódico.

\newpage

% -------------------------------------------------------------------
\subsection*{Iteración 2 (con descuento)}

Ahora $R_1 = [1, 1, 3]$.

\subsubsection*{Paso 1 --- Determinación del valor}

Matriz combinada (Radio en 0,1; Periódico en 2), con $\alpha = 0.9$:

\[
\boxed{\begin{cases}
0.64V_0 - 0.45V_1 - 0.09V_2 = -280 & \text{(D1)}\\[4pt]
-0.09V_0 + 0.37V_1 - 0.18V_2 = -250 & \text{(D2)}\\[4pt]
-0.18V_1 + 0.28V_2 = -300 & \text{(D3)}
\end{cases}}
\]

De (D3): $V_2 = -1071.43 + 0.6429\,V_1$.

En (D2):
\[
-0.09V_0 + 0.37V_1 - 0.18(-1071.43 + 0.6429V_1) = -250
\;\Longrightarrow\; -0.09V_0 + 0.2543V_1 = -442.86
\]
\[
\boxed{V_1 = 0.3540\,V_0 - 1741.49}
\]

Actualizamos $V_2$ con $V_1$:
\[
V_2 = -1071.43 + 0.6429(0.3540V_0 - 1741.49)
     = -1071.43 + 0.2276V_0 - 1119.52
     = -2190.95 + 0.2276\,V_0
\]

En (D1):
\[
\begin{aligned}
0.64V_0 - 0.45(0.3540V_0 - 1741.49) - 0.09(-2190.95 + 0.2276V_0) &= -280 \\
0.64V_0 - 0.1593V_0 + 783.67 + 197.19 - 0.0205V_0 &= -280 \\
0.4602V_0 + 980.86 &= -280 \\
0.4602V_0 &= -1260.86
\end{aligned}
\]

\[
\boxed{V_0 \approx -2739.59}
\]

\[
V_1 = 0.3540(-2739.59) - 1741.49 \approx -969.73 - 1741.49 \approx \boxed{-2711.22}
\]

\[
V_2 = -2190.95 + 0.2276(-2739.59) \approx -2190.95 - 623.60 \approx \boxed{-2814.55}
\]

\[
\boxed{ \boxed{ V_0 = -2739.59,\quad V_1 = -2711.30,\quad V_2 = -2814.37 } }
\]

\textbf{Interpretación:} Los tres valores son más negativos que en la iteración 1, reflejando mayor utilidad descontada. El estado Excelente ($-2814$) pasó de ser el ``menos malo'' al más favorable, porque ahora usa Periódico ($C_{2,3}=-300$) en vez de TV ($C_{2,2}=+130$).

\subsubsection*{Paso 2 --- Mejoramiento}

\textbf{Estado 0:}

\begin{quote}\footnotesize
$k=1$ (Radio):\;
$\alpha\sum P_{0j}V_j = 0.9[0.4(-2739.59) + 0.5(-2711.30) + 0.1(-2814.37)]$ \\
\hspace*{1.8cm} $= 0.9[-1095.8 - 1355.7 - 281.4] = -2459.6$ \\
\hspace*{1.8cm} $\text{valor} = -280 + (-2459.6) = \mathbf{-2739.6}$\quad $\leftarrow$ \textbf{MEJOR}

\vspace{3pt}
$k=2$ (TV):\;
$\alpha\sum P_{0j}V_j = 0.9[0.7(-2739.59) + 0.2(-2711.30) + 0.1(-2814.37)]$ \\
\hspace*{1.8cm} $= 0.9[-1917.7 - 542.3 - 281.4] = -2467.3$ \\
\hspace*{1.8cm} $\text{valor} = -220 + (-2467.3) = -2687.3$

\vspace{3pt}
$k=3$ (Periódico):\;
$\alpha\sum P_{0j}V_j = 0.9[0.2(-2739.59) + 0.5(-2711.30) + 0.3(-2814.37)]$ \\
\hspace*{1.8cm} $= 0.9[-547.9 - 1355.7 - 844.3] = -2473.1$ \\
\hspace*{1.8cm} $\text{valor} = -258 + (-2473.1) = -2731.1$
\end{quote}

\vspace{0.3cm}
\textbf{Estado 1:}

\begin{quote}\footnotesize
$k=1$ (Radio):\;
$\alpha\sum P_{1j}V_j = 0.9[0.1(-2739.59) + 0.7(-2711.30) + 0.2(-2814.37)]$ \\
\hspace*{1.8cm} $= 0.9[-274.0 - 1897.9 - 562.9] = -2461.3$ \\
\hspace*{1.8cm} $\text{valor} = -250 + (-2461.3) = -2711.3$

\vspace{3pt}
$k=2$ (TV):\;
$\alpha\sum P_{1j}V_j = 0.9[0.3(-2739.59) + 0.6(-2711.30) + 0.1(-2814.37)]$ \\
\hspace*{1.8cm} $= 0.9[-821.9 - 1626.8 - 281.4] = -2457.1$ \\
\hspace*{1.8cm} $\text{valor} = -110 + (-2457.1) = -2567.1$

\vspace{3pt}
$k=3$ (Periódico):\;
$\alpha\sum P_{1j}V_j = 0.9[0.0(-2739.59) + 0.7(-2711.30) + 0.3(-2814.37)]$ \\
\hspace*{1.8cm} $= 0.9[0 - 1897.9 - 844.3] = -2468.0$ \\
\hspace*{1.8cm} $\text{valor} = -255 + (-2468.0) = \mathbf{-2723.0}$\quad $\leftarrow$ \textbf{MEJOR}
\end{quote}

\vspace{0.3cm}
\textbf{Estado 2:}

\begin{quote}\footnotesize
$k=1$ (Radio):\;
$\alpha\sum P_{2j}V_j = 0.9[0.1(-2739.59) + 0.2(-2711.30) + 0.7(-2814.37)]$ \\
\hspace*{1.8cm} $= 0.9[-274.0 - 542.3 - 1970.1] = -2507.7$ \\
\hspace*{1.8cm} $\text{valor} = -220 + (-2507.7) = -2727.7$

\vspace{3pt}
$k=2$ (TV):\;
$\alpha\sum P_{2j}V_j = 0.9[0.1(-2739.59) + 0.7(-2711.30) + 0.2(-2814.37)]$ \\
\hspace*{1.8cm} $= 0.9[-274.0 - 1897.9 - 562.9] = -2461.3$ \\
\hspace*{1.8cm} $\text{valor} = +130 + (-2461.3) = -2331.3$

\vspace{3pt}
$k=3$ (Periódico):\;
$\alpha\sum P_{2j}V_j = 0.9[0.0(-2739.59) + 0.2(-2711.30) + 0.8(-2814.37)]$ \\
\hspace*{1.8cm} $= 0.9[0 - 542.3 - 2251.5] = -2514.4$ \\
\hspace*{1.8cm} $\text{valor} = -300 + (-2514.4) = \mathbf{-2814.4}$\quad $\leftarrow$ \textbf{MEJOR}
\end{quote}

\[
\boxed{ R_2 = [1,\; 3,\; 3] = (\text{Radio},\; \text{Periódico},\; \text{Periódico}) }
\]

\textbf{Conclusión:} Ambos métodos (con y sin descuento) convergen en 2 iteraciones a la misma política $[1,3,3]$, confirmando que es la \textbf{política óptima}.

\newpage

% ===================================================================
% G) PLANTEAMIENTO DEL PPL
% ===================================================================
\section{Inciso G --- Planteamiento del PPL}

El PMD puede formularse como un Programa Lineal. Las variables $Y_{ik}$ representan la probabilidad estacionaria conjunta de estar en el estado $i$ y tomar la decisión $k$.

\subsection*{Variables}

9 variables: $Y_{0,1},\, Y_{0,2},\, Y_{0,3},\, Y_{1,1},\, Y_{1,2},\, Y_{1,3},\, Y_{2,1},\, Y_{2,2},\, Y_{2,3}$.

\subsection*{Función objetivo (minimizar)}

\[
\begin{aligned}
\min Z = &\; -280\,Y_{0,1} - 220\,Y_{0,2} - 258\,Y_{0,3} \\
         &\; -250\,Y_{1,1} - 110\,Y_{1,2} - 255\,Y_{1,3} \\
         &\; -220\,Y_{2,1} + 130\,Y_{2,2} - 300\,Y_{2,3}
\end{aligned}
\]

\subsection*{Restricciones}

\textbf{1. Normalización:}
$\displaystyle \sum_{i=0}^{2}\sum_{k=1}^{3} Y_{ik} = 1$

\vspace{4pt}
\textbf{2. Balances de flujo} (para cada $j = 0, 1, 2$):
$\displaystyle \sum_{k=1}^{3} Y_{jk} \;-\; \sum_{i=0}^{2}\sum_{k=1}^{3} Y_{ik} \cdot P_{ij}(k) = 0$

\vspace{4pt}
\textbf{3. No negatividad:} $Y_{ik} \geq 0$ para todo $i,k$.

\subsection*{Estructura}

9 variables $\times$ 4 restricciones. Se resuelve por el \textbf{método Simplex} (técnica de la Gran M).

\subsection*{Solución óptima del PPL}

\begin{center}
\begin{tabular}{c c}
\toprule
\textbf{Variable} & \textbf{Valor óptimo} \\
\midrule
$Y_{1,3}$ & $0.4$ \\
$Y_{2,3}$ & $0.6$ \\
Resto     & $0$ \\
\bottomrule
\end{tabular}
\end{center}

\textbf{Valor óptimo:} $Z^* = -282.00$ \quad $\Rightarrow$ \quad \textbf{Utilidad = \$282.00 / semana}.

\subsection*{Política determinística}

Se calcula $D_{ik} = Y_{ik} / \sum_{k} Y_{ik}$:

\begin{center}
\begin{tabular}{c|c|c|c|c}
\toprule
\textbf{Estado $i$} & $D_{i,1}$ (Radio) & $D_{i,2}$ (TV) & $D_{i,3}$ (Per.) & \textbf{Decisión} \\
\midrule
0 (Regular)   & -- & -- & -- & 1 (Radio)* \\
1 (Bueno)     & 0 & 0 & \textbf{1.0} & \textbf{3 (Periódico)} \\
2 (Excelente) & 0 & 0 & \textbf{1.0} & \textbf{3 (Periódico)} \\
\bottomrule
\end{tabular}
\end{center}

\begin{quote}
\footnotesize *El estado 0 tiene $Y_{0k}=0$ para las tres decisiones porque es \textbf{transitorio} bajo la política óptima (nunca se regresa a él). Cualquier decisión es equivalente; se asigna $k=1$ (Radio) por convención.
\end{quote}

\[
\boxed{ R_{\text{PL}} = [1,\; 3,\; 3] = (\text{Radio},\; \text{Periódico},\; \text{Periódico}) }
\]

\vspace{0.5cm}
\begin{center}
\rule{0.7\textwidth}{0.4pt}\\[0.3cm]
{\large \textbf{Conclusión final}}\\[0.3cm]
Los tres métodos ---mejoramiento sin descuento, con descuento ($\alpha=0.9$), y PL--- \textbf{convergen a la misma política óptima:}

\[
\fbox{\Large $R^* = [\text{Radio},\; \text{Periódico},\; \text{Periódico}]$}
\]

con \textbf{utilidad esperada de \$282.00 por semana}.
\end{center}

\end{document}
